\chapter*{Conclusion}

\lettrine{A}u terme de cette étude menée au sein du Département du patrimoine, de la conservation et de la numérisation de la Bibliothèque interuniversitaire Cujas, force est de constater que la modernisation de l'accès au patrimoine juridique ne saurait se réduire à une simple opération de migration informatique. Si l'obsolescence de l'infrastructure historique sous XTF et la nécessité d'évaluer le nouveau prototype Codex ont constitué le point de départ de l'étude, afin de saisir l'ensemble des enjeux nécessaire à cette refonte, cela nous a rapidement amené à dépasser le cadre purement technique du logiciel et commment celui-ci doit intéragir avec les autres. En effet, s'afffranchir d'une dette technique accumulée au fil des années et rompre l'isolement documentaire d'une institution de référence requiert bien plus que l'implémentation de nouveaux standards du Web sémantique. Cette refonte amène le besoin de repenser en profondeur l'articulation entre les pratiques professionnelles des équipes et les usages d'une communauté de chercheurs en droit aux exigences spécifiques. Ces quelques mois de stage ont ainsi amené à réfléchir à l'objectif principal de cette refonte, d'identifier clairement le public qui utilisent aujourd'hui la bibliothèque numérique ainsi que les attentes aussi bien du côté métier que du côté usagers. Ce travail a un impact concret pour les agents de la bibliothèque Cujas qui ont une vision plus claire des attentes et des enjeux liés à cette nouvelle plateforme, ainsi que les futurs chantiers a mené avant de pouvoir entamer cette refonte et de rendre leur plateforme interopérable et comptable avec les autres bibliothèques numériques.

En définissant des personae et en formalisant leurs parcours de recherche, les équipes ont acquis une vision bien plus précise des attentes fonctionnelles. Ce projet a également ouvert un espace de dialogue fructueux entre le Service du patrimoine numérisé (SPN), le département informatique et la communauté scientifique. Il a ainsi posé les bases d'un langage commun et d'une culture projet partagée, deux prérequis essentiels avant d'engager les développements futurs.

CujasNum est l'une des premières bibliothèques numériques à avoir été développé afin de répondre précisément aux attentes des chercheurs en histoire du droit. Toutefois, les pratiques de recherche contemporaines, il est du rôle de la bibliothèque Cujas de faire évoluer sa bibliothèque numérique afin d'asseoir sa légitimité pour devenir une bibliothèque de référence juridique. Cette démarche a mis en lumière une tension fondamentale entre d'un côté l'attachement fort de la communauté des chercheurs à des méthodes de consultation traditionnelles et de l'autre côté qu'il est primordial pour la bibliothèque de s'affranchir de son isolement technique.

Pour les chercheurs, l'outil numérique doit avant tout recréer la manipulabilité et la logique d'un livre papier, une exigence qui se traduit par une forte demande pour des tables des matières hiérarchiées et une océrisation exhaustive des textes. Pour la bibliothèque, il est devenu primordial pour elle de s'affranchir de l'isolement dans lequel l'a enfermé CujasNum. L'institution se doit en effet de répondre aux impératifs contemporains d'interopérabilité, d'alignement des métadonnées sur des référentiels nationaux et d'intégration au Web sémantique.

Le coeur de la démarche adoptée ne prétend pas l'exhaustivité, mais de poser les fondamentations sur lesquelles se reposer afin de continuer la réflexion entamée. Cette étude a donc permis d'interroger ce paradoxe : comment inscrire les collections de Cujas dans des réseaux documentaires interconnectés et ouverts, tout en concevant, à travers l'étude du prototype Codex, des services sur-mesure capables de respecter l'ergonomie et la spécificité de la recherche en histoire du droit? L'expérience de stage montre la réinformatisation d'une institution patrimoniale n'est pas seulement une question de ressources humaines, financiés et organisationnelles, mais engage à la bibliothèque dans un chantier au long cours. 

Bien que les pistes tracées offrent une feuille de route fonctionnelle solide pour l'établissement, il convient de les évaluer en tenant compte de la durée limitée de ce stage. Ces quatre mois de stage ont permis d'établir une étude en amont solide, mais ne peut pas être considéré comme une fois en soi et nécessite une poursuite de cet énorme projet. Plusieurs points importants ont dû être écartés de notre périmètre immédiat, à l'instar du traitement spécifique des périodiques et des archives — dont les structures de métadonnées diffèrent grandement des monographies —, ou encore du nettoyage et de l'alignement de masse des métadonnées existantes. Par ailleurs, des questions émergentes comme l'application de l'intelligence artificielle pour l'indexation sémantique automatique ou la correction automatisée de l'OCR défectueux restent des territoires d'avenir à explorer. Un autre point écarté fut la question de l'anonymisation des noms de parties qui est devenue obligatoire lorsque l'on met en ligne des revues de jurisprudence, nécessitant une étude à part sur cette question.
Néanmoins, sur le plan de l'infrastructure, l'avenir de la plateforme passe désormais par des choix technologiques clairs. Le développement souverain et sur-mesure sous Python-Django engagé avec le prototype Codex prouve sa pertinence pour s'affranchir des solutions clés en main du marché, mais il demandera à l'institution de pérenniser et de valoriser les compétences de développement technique en interne. De plus, le déploiement du protocole IIIF (International Image Interoperability Framework) et de sa visionneuse Mirador devra être poussé au-delà de la simple visualisation : il s'agira de permettre aux chercheurs de manipuler, d'annoter et de comparer directement des sources juridiques dispersées dans d'autres bibliothèques numériques mondiales, transformant CujasNum en un nœud de recherche interactif de premier plan.

En dernière analyse, la refonte de CujasNum dépasse largement les seuls enjeux d'architecture logicielle ou de transition bibliographique pour toucher à la mission politique, sociale et démocratique de la bibliothèque universitaire. En s'engageant résolument dans la voie de l'interopérabilité et en s'ouvrant aux principes directeurs de la Science Ouverte, la bibliothèque Cujas ne modernise pas simplement un outil informatique : elle réaffirme son rôle de passeur de savoirs. C'est en devenant un espace d'interconnexion vivant au sein de l'Enseignement supérieur et de la Recherche que l'établissement continuera de faire vivre son patrimoine d'exception. En garantissant un accès sémantique, fluide et universel aux sources historiques du droit, la bibliothèque Cujas permet à chacun — du doctorant le plus spécialisé au citoyen désireux d'étudier les fondements de nos lois — de se réapproprier cette mémoire commune, assurant ainsi le rayonnement durable de la culture juridique française.